\documentclass[11pt]{article}
\usepackage[margin=1in]{geometry}
\usepackage{amsmath,amssymb}
\usepackage{hyperref}
\PassOptionsToPackage{hyphens}{url}
\usepackage{url}

\title{An independent machine check of a claimed Riemann Hypothesis proof}
\author{Dr Alex (Aco) Kovaceski\\AxoQuant (Kovacorp Pty Ltd), Canberra, ACT\\\texttt{alex@axoquant.com}}
\date{26 November 2026}

\begin{document}
\maketitle

\begin{abstract}
A paper posted at \url{https://alpo.ge/remark.pdf} claims a proof of the Riemann Hypothesis, developed with Claude and formalized in Lean without assumptions. The only public evidence is a Lean archive of 175,187 files and 17.4 million lines. This report records an independent check of that claim. The archive was statically screened, compiled from source with the pinned toolchain and mathlib revision across a Slurm cluster, and the four final theorems were examined for axiom dependencies. The final theorems depend only on the standard axioms of Lean's logic (propext, Classical.choice, Quot.sound). The claim is verified. The receipt and reproduction commands are published alongside this note. The result is independently corroborated by Michael Rubinstein (University of Waterloo), including a replay through Lean's kernel and a check by the independent nanoda kernel.
\end{abstract}

\section{The claim}
The paper \emph{remark.pdf} (L. Alp\"oge) states that a proof of the Riemann Hypothesis has been discovered with Claude, that the arguments are formalized in Lean without assumptions, and that the method resolves the Generalized Riemann Hypothesis for Dirichlet L-functions. The paper declines to publish the human-readable argument. The only public artifact is the Lean archive, linked as fourteen tarball parts at \url{https://alpo.ge}.

The question this report answers is narrow: does the archive actually prove the Riemann Hypothesis? Not whether the argument is understood, or whether it is elegant, or whether it is the shortest. Whether it compiles, and whether the final theorems depend on anything beyond Lean's standard logical axioms.

\section{The artifact}
The archive is a Lean 4 project pinned to \texttt{leanprover/lean4:v4.33.1} and mathlib revision \texttt{0df444a360eaa60ab8c11dca51a86af692955474}. It contains 175,187 \texttt{.lean} files, 17,374,215 lines, and 1.03 GB of source. For scale, mathlib itself is about 7,000 files. The final theorems live in \texttt{src/Main.lean}:
\begin{verbatim}
theorem riemannHypothesis : RiemannHypothesis := rhOfAll
theorem riemannZeta_eq_zero_re : RiemannZetaZerosOnLine := rhOfAll
theorem LFunction_eq_zero_re : DirichletLZerosOnLine := claimAll.2.1
theorem LFunction_eq_zero_re_of_isPrimitive : PrimitiveDirichletLZerosOnLine := critAll
\end{verbatim}

The first is the Riemann Hypothesis as stated in mathlib. The other three are the corresponding statements for Dirichlet L-functions: nonvanishing right of $1/2$, zeros on the line, and the primitive-character form.

\section{Static checks}
Five checks were run before the build, each able to fail.

\textbf{Statement integrity.} \texttt{RiemannHypothesis} is mathlib's canonical definition at the pinned revision: every nontrivial zero of the Riemann zeta function has real part $1/2$. No shadow definition of \texttt{riemannZeta} exists in the archive, so the premise cannot be vacuously false.

\textbf{Unsoundness sweep.} Zero occurrences of \texttt{axiom}, \texttt{constant}, \texttt{sorry}, \texttt{admit}, \texttt{unsafe}, \texttt{implemented\_by} or \texttt{opaque} across all 17.4 million lines. No axiom-injecting elaborator.

\textbf{Vacuity sweep.} The load-bearing statements (\texttt{EveryCharacter}, \texttt{Row0b}, \texttt{TLiftK}, \texttt{HalfPlane}, \texttt{Claim.Zeta}) hold real mathematical content; none is defined as \texttt{True}. \texttt{Claim.Zeta} is definitionally mathlib's \texttt{RiemannHypothesis}, shown by \texttt{Iff.rfl}.

\textbf{Wiring integrity.} 753,569 internal imports, zero unresolved.

\textbf{Duplication.} 14 duplicate files out of 175,187.

None of these settles the claim. Elaboration fails in ways no static scan sees. The compiler is the instrument.

\section{The build}
The archive was built from source with the pinned toolchain and the pinned mathlib revision. The build was sharded across the AxoQuant Slurm cluster, with the closing pass and the receipt run on the coordinator. The build record, including the operational failures and their causes, is published with the review. The failures were all infrastructure (disk exhaustion from intermediate dumps, a flag the toolchain does not accept, and a cleanup that deleted directories mid-build), not elaboration errors in the archive.

\section{The receipt}
After the full archive compiled, \texttt{lake env lean GovernorCheck.lean} produced, on 24 November 2026:
\begin{verbatim}
'riemannHypothesis' depends on axioms: [propext, Classical.choice, Quot.sound]
'riemannZeta_eq_zero_re' depends on axioms: [propext, Classical.choice, Quot.sound]
'LFunction_eq_zero_re' depends on axioms: [propext, Classical.choice, Quot.sound]
'LFunction_eq_zero_re_of_isPrimitive' depends on axioms: [propext, Classical.choice, Quot.sound]
\end{verbatim}
The three axioms listed are the standard axioms of Lean's underlying logic: propositional extensionality, the axiom of choice, and quotient soundness. They are present in every mathlib proof. No custom axioms, no \texttt{sorry}, no interface assumptions.

\section{Verdict}
Verified. The archive compiles in full with the pinned toolchain and mathlib revision, and the four final theorems depend only on the standard logical axioms. The final theorem is mathlib's canonical statement of the Riemann Hypothesis. A compiled, axiom-clean term of that type is a proof of the Riemann Hypothesis, and the same holds for the three GRH forms. The claim is demonstrated, not just asserted.

\section{Independent corroboration}
Michael Rubinstein (University of Waterloo) reports an independent verification of the archive on 21 November 2026 at \href{https://www.math.uwaterloo.ca/~mrubinst/alpoge_claude_grh/}{the Waterloo verification page}: all 175,187 modules compile; the five theorems depend only on Lean's three standard axioms; a replay of every module through Lean's kernel accepted all 175,187; and a second, independently written kernel (nanoda) checked the 1,484,729 declarations the theorems depend on, with no typechecker errors. Further independent checks are reported by Andrew Sutherland (MIT), Elliot Arledge, and a status dashboard maintained by Bartosz Naskr\k{e}cki. The Lean and Mathlib communities' joint statement of 21 November 2026, endorsed by the Lean FRO and the Mathlib Initiative, certifies that both the statement and the Lean proof of the Riemann Hypothesis are correct. Rubinstein also publishes a 265-page human-readable account of the argument for the zeta function (version 2, 22 November 2026), reducing the Riemann hypothesis to a main inequality and thirteen numbered hypotheses.

\section{A clarification from the author}
On 23 November 2026 Alp\"oge posted that the 250,000-page figure in his footnote was a joke (\href{https://www.businessinsider.com/anthropic-levent-alpoe-riemann-hypothesis-page-length-joke-angers-mathematicians-2026-11}{Business Insider, 24 November 2026}). The Lean archive remains the only formal evidence, and it is what was checked. The page-length episode does not bear on the verification.

\section{Reproduction}
\begin{verbatim}
# toolchain
curl -sSf https://raw.githubusercontent.com/leanprover/elan/master/elan-init.sh | sh -s -- -y
elan toolchain install leanprover/lean4:v4.33.1
# source archive: fetch lean.tar.gz.part01 .. part14 from https://alpo.ge,
# cat into lean.tar.gz, extract
cd lean && lake build Main
lake env lean GovernorCheck.lean
\end{verbatim}
The check file, the review notes, the shard tooling and the build record are at \url{https://github.com/axoquant/rh-governance-review}. The citable record with a DOI is at \url{https://zenodo.org/records/24519677}.

\section{Limitations}
This report certifies that the archive compiles and that the final theorems are axiom-clean. It does not certify that the proof is the simplest, that the mathematics is understood, or that the human-readable argument exists. Those are separate questions. An independent outline of the argument by S. Slijep\v{c}evi\'c (\url{https://cantabpi.com/grh.pdf}) records what a reader could and could not understand of the development.

\end{document}
